% TOP_PROBLEM: {"id": 2618, "title": "Kourovka Notebook Problem 21.109", "category_id": 17, "category": {"id": 17, "name": "group_theory", "display_name": "Group Theory", "description": "Problems about groups, group actions, representations, and related algebraic structures.", "slug": "group-theory", "order_index": 17, "created_at": "2026-07-31T15:26:25.671Z"}, "status": "open", "problem_number": "KOU-21.109", "set_id": 10}
% FIRST_POSTED: 2026-10-01
% ENTRY_KIND: statement_only
% UnsolvedMath ID 2618; source catalogue code KOU-21.109
% !TeX program = lualatex
% Definition and sources only; this file is not an LLM solution attempt.
\documentclass[11pt,a4paper]{article}
\usepackage[margin=25mm]{geometry}
\usepackage{fontspec}
\setmainfont{lmroman10-regular.otf}[BoldFont=lmroman10-bold.otf,
 ItalicFont=lmroman10-italic.otf,BoldItalicFont=lmroman10-bolditalic.otf]
\usepackage{amsmath,amssymb,mathtools}
\usepackage{unicode-math}
\setmathfont{latinmodern-math.otf}
\AtBeginDocument{\let\setminus\smallsetminus}
\usepackage{xurl,hyperref}
\hypersetup{colorlinks=true,urlcolor=blue}
\setcounter{secnumdepth}{0}
\setlength{\parindent}{0pt}
\setlength{\parskip}{5pt}
\setlength{\emergencystretch}{3em}
\begin{document}
\section*{Kourovka Notebook Problem 21.109}\label{2618}
{\small UnsolvedMath ID 2618 · KOU-21.109 · Group Theory · Kourovka Notebook - New Problems, Issue 21}
\subsection{Definitions and mathematical statement}
For a group $G$, put $G^{(0)}=G$ and define the derived series recursively by $G^{(j+1)}=[G^{(j)},G^{(j)}]$, where $[H,H]$ is generated by $a^{-1}b^{-1}ab$ for $a,b\in H$. The group is solvable if this series reaches the trivial group. Its derived length $\operatorname{dl}(G)$ is the least $j\ge0$ with $G^{(j)}=1$; in particular $\operatorname{dl}(1)=0$.

For a finite group and an irreducible complex character $\chi$, let $\chi(1)$ be its representation dimension and $\ker\chi$ its representation kernel. Define
\[
 \operatorname{cod}(\chi)=\frac{[G:\ker\chi]}{\chi(1)},\qquad
 \operatorname{Cod}(G)=\{\operatorname{cod}(\chi):\chi\in\operatorname{Irr}(G)\}.
\]
Codegrees are positive integers. The braces form a set: equal codegrees from different characters are counted only once.

The conjecture is that every finite solvable group satisfies
\[
 \operatorname{dl}(G)\le|\operatorname{Cod}(G)|-1.
\]
Thus the number of successive nontrivial steps in the derived series should be controlled by the number of distinct irreducible character codegrees. The principal character contributes codegree one in all groups; this explains the subtraction in the proposed bound. The assertion includes the trivial group, for which both sides are zero. The target concerns derived length, and the definition does not replace it by the length of another normal series.
\subsection{Short English statement}
Is a finite solvable group’s derived length at most one less than its number of distinct irreducible character codegrees?
\subsection{Sources}
\begin{itemize}
\item[S1] ULAM.AI and the UnsolvedMath Contributors, supplied problems.json, record 2618 (KOU-21.109), Kourovka Notebook - New Problems, Issue 21. Catalogue identity and source transcription; CC BY 4.0.
\url{https://huggingface.co/datasets/ulamai/UnsolvedMath}.
\item[S2] The Kourovka Notebook, 21st edition, version 47 (30 September 2026), new problems, Problem 21.109. Expanded formulation of the supplied catalogue statement.
\url{https://arxiv.org/pdf/1401.0300v47}.
\end{itemize}
\end{document}
